\section{Background and Related Work}
\label{sec:background}

\subsection{Invariant Subspaces of Bounded Operators}

Let $T$ be a bounded linear operator on a separable Hilbert space $H$.
A closed subspace $V \subseteq H$ is \emph{invariant} under $T$ if
$T(V) \subseteq V$. The \emph{invariant subspace problem} asks whether
every such $T$ on an infinite-dimensional separable $H$ admits a
non-trivial closed invariant subspace---one other than $\{0\}$ or $H$
itself. The question is open in the Hilbert space case; in finite
dimensions every operator trivially admits invariant subspaces (its
eigenspaces, in particular). Two structural facts about the literature
matter here.

First, when an operator admits an invariant subspace, its action
decomposes: writing $T$ in block form with respect to the splitting
$H = V \oplus V^\perp$ shows that the upper-right block vanishes and
the operator's geometry becomes legible through the geometry of its
restrictions to invariant subspaces. The existence of invariant
subspaces is the existence of a useful coordinate system for the
operator.

Second, the affirmative results form a hierarchy. Aronszajn and Smith
\cite{aronszajn1954invariant} established that completely continuous
(compact) operators always admit non-trivial invariant subspaces.
Lomonosov \cite{lomonosov1973invariant} extended this dramatically: any
operator commuting with a non-zero compact operator admits a
non-trivial invariant subspace. Read \cite{read1984solution}
constructed counter-examples on $\ell^1$, and Argyros and Haydon
\cite{argyros2011hereditarily} produced a Banach space on which every
bounded operator has the form scalar-plus-compact (so every operator
has invariant subspaces, but the construction requires extreme
hereditary indecomposability). The Hilbert space case remains open.
Beauzamy's monograph \cite{beauzamy1988introduction} gives the standard
treatment.

Transformer forward passes are not bounded linear operators. They are
nonlinear maps composed of linear projections, multi-head attention
(itself a nonlinear function of the input through softmax), gated
feed-forward layers, and residual connections. The classical theorems
do not apply. We use the language of invariant subspaces in an
\emph{approximate, empirical} sense: a $k$-dimensional subspace
$V \subseteq \R^d$ is approximately $\epsilon$-invariant under a map
$f: \R^d \to \R^d$ on a sample $X$ if
$\fnorm{(I - P_V) f(P_V X)} / \fnorm{f(P_V X)} \leq \epsilon$, where
$P_V$ is the orthogonal projection onto $V$. This is a sample-level
operationalization of the structural intuition the classical theory
formalizes; we make no claim about the limiting infinite-dimensional
case.

\subsection{Mechanistic Interpretability: Circuits and Sparse Autoencoders}

The mechanistic interpretability program seeks to identify the
computational primitives implemented by trained neural networks,
typically by reverse-engineering specific circuits or by decomposing
the residual stream into interpretable components.

The circuits-level work has identified increasingly specific mechanisms.
Olah et al.~\cite{olah2020zoom} articulated the program through case
studies in vision models. Elhage et al.~\cite{elhage2021mathematical}
provided a mathematical framework for transformer circuits.
Olsson et al.~\cite{olsson2022context} identified \emph{induction
heads} as a specific attention-head pattern responsible for the
majority of in-context learning behavior in large transformers.
Wang et al.~\cite{wang2023interpretability} reverse-engineered a
multi-component circuit for indirect-object identification in GPT-2
Small. Geva et al.~\cite{geva2021transformer} showed that transformer
feed-forward layers behave as key-value memories. Meng et
al.~\cite{meng2022rome} demonstrated that specific factual associations
can be located and edited in MLP weights. Conmy et
al.~\cite{conmy2023towards} introduced automated circuit discovery.

The polysemanticity problem---that single neurons typically respond to
multiple unrelated concepts---was formalized by Elhage et
al.~\cite{elhage2022superposition} as the \emph{superposition}
hypothesis: networks pack more features than they have neurons by
representing them as sparse linear combinations in a higher-dimensional
implicit feature space. The natural inversion of this hypothesis is to
\emph{recover} the implicit feature space via sparse dictionary
learning, which is what sparse autoencoders do.

The SAE program has scaled rapidly. Bricken et
al.~\cite{bricken2023monosemanticity} demonstrated that sparse
autoencoders trained on the residual stream of small models recover
features that are individually interpretable across many examples.
Cunningham et al.~\cite{cunningham2024sparse} extended this to larger
models with substantially larger feature dictionaries.
Templeton et al.~\cite{templeton2024scaling} performed the procedure on
Claude 3 Sonnet, recovering on the order of millions of features,
including features that respond to abstract concepts (deception, code
errors, the Golden Gate Bridge) across modalities. Gao et
al.~\cite{gao2024scaling} demonstrated systematic scaling of SAE
training and evaluation. Lieberum et al.~\cite{lieberum2024gemma}
released open SAE checkpoints for the Gemma 2 family. Marks et
al.~\cite{marks2024sparse} showed that SAE features can be assembled
into causal circuits at the feature level. The empirical case for SAE
features as the right basis for the residual stream is now substantial
across multiple research groups, multiple model families, and multiple
scales.

The universality hypothesis---that different networks recover similar
features---has accumulated specific evidence. Chughtai et
al.~\cite{chughtai2023toy} showed that small networks trained on group
operations recover the same algorithmic decompositions across random
seeds and architectural variations. Gurnee et al.~\cite{gurnee2024universal}
identified \emph{universal neurons} in GPT-2 models that activate on the
same patterns across independently trained models. Nanda et
al.~\cite{nanda2023progress} showed that grokking dynamics correspond
to discrete phase transitions in the recovered circuit, with similar
dynamics across initializations.

\subsection{Cross-Model Representational Alignment}

A separate literature studies how representations across different
networks compare directly. The methodological starting point is the
observation that comparing two representation matrices requires
choosing which symmetries one wants the comparison to be invariant to.

Word-embedding alignment is the historical precedent. Mikolov et
al.~\cite{mikolov2013distributed} showed that word embeddings trained
independently on different corpora have approximately the same internal
structure. Smith et al.~\cite{smith2017offline} and Conneau et
al.~\cite{conneau2017word} demonstrated that bilingual lexicons could
be recovered without parallel data by orthogonal Procrustes alignment
of monolingual word embedding spaces. The success of this procedure
across language pairs and across embedding methods is the strongest
empirical case for substrate-independent linguistic structure that
exists in the static-embedding literature.

For learned-feature representations of the kind transformers produce,
the methodological landscape is broader. Raghu et al.~\cite{raghu2017svcca}
introduced singular vector canonical correlation analysis (SVCCA) as a
representation-similarity measure. Morcos et al.~\cite{morcos2018insights}
developed projection-weighted CCA (PWCCA) addressing some of SVCCA's
sensitivities. Kornblith et al.~\cite{kornblith2019similarity}
introduced centered kernel alignment (CKA) as a measure invariant to
orthogonal transformations and isotropic scaling but not to invertible
linear transformations, and showed it discriminates between models in
ways that the previous measures do not. Williams et
al.~\cite{williams2021generalized} proposed a family of generalized
shape metrics including CKA and Procrustes as special cases. Klabunde
et al.~\cite{klabunde2024similarity} provide a recent survey of the
methodological landscape.

Bansal et al.~\cite{bansal2021revisiting} introduced model stitching as
a complementary tool: rather than measuring representational similarity
directly, train a thin connecting layer between the lower half of one
model and the upper half of another, and observe whether the stitched
model retains performance. They argue that model stitching reveals
representational compatibilities that scalar similarity measures miss.
Lenc and Vedaldi \cite{lenc2015understanding} are an earlier
representative of the same methodological intuition applied to
convolutional networks.

The most recent and most directly related synthesis is Huh et
al.~\cite{huh2024platonic}, the \emph{Platonic Representation
Hypothesis}: representations across models are converging, with scale
and capability, toward a shared statistical model of reality. Their
evidence assembles results from many of the works cited above into a
single hypothesis. The COSI program of the present paper is a specific
operationalization of one consequence of their hypothesis: if
representations converge, they should admit a sharp orthogonal
alignment, with a residual significantly below the chance level set by
the appropriate null distribution. The contribution we add is the
operator-theoretic framing that explains \emph{why} the convergence
should produce alignable subspaces, not merely correlated activations.

\subsection{Substrate-Independence: Cognitive Science, Cybernetics, Free-Energy}

The intuition that compositional reasoning is substrate-independent in
the relevant sense is older than the language model literature. Marr's
three levels of analysis \cite{marr1982vision} distinguish the
computational, the algorithmic, and the implementational levels and
argue that the computational level is independent of the substrate
that performs it. Ashby's \emph{law of requisite variety}
\cite{ashby1956introduction} establishes a structural constraint:
any controller must contain at least as much variety as the disturbance
it regulates, regardless of the controller's substrate. Wiener's
cybernetics \cite{wiener1948cybernetics} provides the framework within
which substrate-independent feedback structure becomes the object of
study. Maturana and Varela's autopoiesis
\cite{maturana1980autopoiesis} extends the framework to
self-sustaining computational systems whose organization is preserved
even as their components turn over.

Friston's free-energy principle \cite{friston2010free} provides a
unifying mathematical framework: any self-maintaining system that
resists dispersion into thermal equilibrium must minimize a variational
free energy that bounds surprise on its sensory inputs. The principle
predicts that any such system---biological or artificial---should
develop internal states that approximate the posterior over its
environment's hidden causes, and that these internal states should be
substrate-independent in the relevant sense, since the variational
objective is. Achille and Soatto \cite{achille2018emergence} derive an
analogous result specifically for deep networks under information
bottleneck regularization \cite{tishby2015deep}, showing that invariance
and disentanglement emerge as consequences of an information-theoretic
objective rather than as architectural choices.

Computational neuroscience has accumulated direct evidence for
substrate-independent invariants in biological computation. Place cells
\cite{okeefe1971hippocampus} and grid cells
\cite{hafting2005microstructure} implement spatial maps with consistent
geometric structure across individual animals. Yamins and DiCarlo
\cite{yamins2016using} demonstrated that goal-driven deep networks
develop intermediate representations that linearly predict primate
visual cortex responses. Gallego et al.~\cite{gallego2017neural} and
Saxena and Cunningham \cite{saxena2019towards} articulate a \emph{neural
population doctrine}: the natural unit of analysis is the population's
collective state in a low-dimensional manifold, not the individual
neuron, and the manifolds have stable geometric structure across
animals and behaviors. The convergence of these observations with the
universality findings in artificial networks is, in our reading, what
the operator-theoretic framework would predict if the framework
captures the actual structure of the computation.

\subsection{Sequence-Model Lineage}

Briefly: the architecture under study (transformers) sits in a lineage
of sequence models including LSTMs \cite{hochreiter1997long}, structured
state-space models (S4 \cite{gu2022s4}, Mamba \cite{gu2024mamba}), and
more recent hybrid architectures combining attention with state-space
components. The Qwen3-Next family used in our Phase~1 is an instance of
this hybrid design. The operator-theoretic framing predicts that the
isomorphism, if it holds, should hold across these architectural classes
and not only within them. We address this directly in the Phase~1
specification (Section~\ref{sec:discussion}). The transformer
architecture itself \cite{vaswani2017attention} is treated as
representative of the dense pure-attention class; we make no claims
specific to its design choices that would not also apply to any
architecture implementing approximate-finite-state computation over
sequence inputs.
